\section{Notation and implementation conventions}
\label{app:notation}

\begin{longtable}{@{}p{0.25\textwidth}p{0.67\textwidth}@{}}
\toprule
Symbol & Meaning \\
\midrule
\endhead
$\PR$ & Primitive-recursive, always for a fixed definition unless a bounded program table is explicitly supplied \\
$\B$ & Numerical Booleans $\{0,1\}$ \\
$u,x,p$ & Paired numerical input/certificate/representation codes \\
$k$ & Requested output precision \\
$h,r$ & Interaction horizon or remaining resource \\
$s$ & Encoded deterministic or consumer state \\
$\Viable(x,r,s)$ & Backward Boolean saying a successful finite continuation exists \\
$\Local(x,r,s,a)$ & Local admissibility of answer $a$ \\
$E$ & Observational equivalence, not necessarily literal equality \\
$\delta$ & Diameter/error bound for stable relational outputs \\
$\Prf\ \PRAAx\ \Gamma\ \varphi$ & Object-level derivability in ordinary iPRA \\
$\NPrf$ & Derivability in the native scalar calculus \\
\bottomrule
\end{longtable}

All finite sequences and graphs may be encoded by any standard primitive-recursive pairing/list scheme with explicit bounds.  Proofs use only closure properties and bounded lookup; no specific optimal coding is mathematically essential.  Implementations should nevertheless share one canonical repository library to avoid repeated conversion proofs.

Rational inequalities are intended to be reduced to signed-integer cross products with positive denominators.  Every use of an order decision is therefore a decision about exact integers, not an arbitrary represented real.
